\section{서론}
\label{sec:introduction}

디지털 제품 여권(Digital Product Passport, DPP)은 제품의
소재·환경정보·이력을 제품에 연결하여 여러 주체가 활용할 수 있게
하는 체계다. EU 집행위원회는 이를 통해 소비자의 구매 결정,
수리·재활용과 관계 기관의 규정 준수 점검을 지원하고자
한다.~\cite{ec-dpp-overview}
재활용 원료 함량과 탄소배출 정보는 DPP에서 다루는
지속가능성 정보의 예다.~\cite{udokwu-crass-2026}
그중 원료 관련 주장을 뒷받침하는 계산·배분 정보는 최초 증빙과
이후의 가공·이전 내역에 함께 의존한다. 예를 들어 ISCC PLUS의
원료 이력관리 기준은 인증 원료의 입출고량과 배분 내역을
기록하도록 한다.~\cite{iscc-plus-coc-2025}
이러한 배분·이전 결과는 다음 기업의 계산 근거로 다시 쓰이지만,
그 정보를 이용하는 기업이나 검사기관은 앞선 거래에 직접
참여하지 않았을 수 있다. 따라서 표시된 값만 받아보는 데서
나아가, 그 값이 정해진 계산·배분 규칙에 따라 작성되었는지도
확인할 수단이 필요하다.

이러한 확인은 신뢰하는 기관이 자료를 검사해 결과를 알려주는
방식으로도 가능하다. 다만 결과만 전달받는 참여자는 계산·배분의
적절성을 그 기관의 판단에 의존해야 한다. 본 논문은 계산·배분에
관한 이러한 의존을 줄이기 위해, 거래 당사자가 다른 참여자도 직접
확인할 수 있는 검증 자료를 제공하는 방식을 택한다.
개별 거래의 검증 결과는 이후 거래를 판단하는 근거가 되므로,
그 결과가 반영된 승인 이력도 함께 관리해야 한다.
참여자마다 서로 다른 승인 이력을 기준으로
삼으면 유효하다고 보는 거래도 달라질 수 있기 때문이다.
이에 여러 참여자가 같은 승인 이력과 원장 상태를 기준으로
검증하고 결과를 관리하도록 하며, 이를 지원하는 방법 중 하나로
블록체인 원장을 사용한다. 블록체인은 거래 이력을 여러
참여자가 공동으로 유지하고, 확정된 기록을 임의로 변경하기
어렵게 한다.~\cite{nist-blockchain-2018}
이 과정에서 검증에 필요한 자료는 원장 참여자에게 공유되어
이후에도 조회할 수 있는 이력으로 남는다.

그런데 직접 검증할 근거를 제공하려고 원료별 투입량이나
거래 상대를 그대로 공개하면, 제품 배합과 조달 관계까지
드러날 수 있다.~\cite{udokwu-crass-2026}
따라서 검증에 필요한 근거는 제공하면서 그 근거를 만드는 데
쓰인 상세 원자료는 보호해야 한다.
\emph{영지식증명}은 비밀 자료를 공개하지 않으면서도 그 자료가
정해진 조건을 만족한다는 사실을 다른 사람이 확인하게 하는
암호학적 방법이다.~\cite[20.3.5절]{boneh-shoup-2020}
이를 이용하면 원료의 양과 실제 거래 연결을 숨긴 상태에서도
계산·이전 규칙의 준수를 검증할 수 있다.

이처럼 거래가 규칙에 맞게 처리되었음을 확인하더라도,
그 근거가 된 증빙의 오류가 나중에 밝혀질 수 있다.
예를 들어 재활용 원료라는 증빙을 바탕으로 원료를 두 기업에
배분하고, 각 기업이 이를 가공·이전한 뒤에야 최초 증빙이
잘못되었음을 발견할 수 있다. 각 거래의 계산이 규칙에
맞았더라도, 그 결과로 작성된 원료 정보는 잘못된 최초 증빙에
의존한다. 이 정보가 다음 거래에 다시 쓰이면 오류의 영향도
새로운 거래로 이어질 수 있다. 따라서 문제가 된 증빙에서
파생된 후속 정보를 찾아, 그 정보가 새로운 거래의 근거로
계속 쓰이는 것을 제한할 방법이 필요하다.

이러한 후속 정보를 찾으려면 어떤 거래의 결과가 다음 거래의
근거로 쓰였는지 확인할 수 있어야 한다. 그러나 이 조사에
필요한 원료의 배분·이전 관계는 앞서 기업정보 보호를 위해
숨기고자 한 정보이기도 하다. 따라서 일반 참여자에게는
실제 거래 연결을 공개하지 않으면서도, 승인된 감사자는
그 연결을 따라 후속 정보를 찾을 수 있어야 한다.
거래마다 이 연결을 암호화해 남겨 두면, 승인된 감사자가
이를 복원하여 후속 이력을 조사할 수 있다.

암호화해 남긴 정보가 실제 거래와 맞는지는 또 다른 문제다.
계산·이전 규칙만 검사하는 증명은 별도로 남긴 감사 자료의
정확성까지 자동으로 검사하지 않는다. 그러므로 거래 계산이
맞더라도 실제 사용한 기록 대신 다른 기록을 암호화해 남기면,
감사자는 잘못된 연결을 따라갈 수 있다. 대상을 올바르게
찾더라도 감사자가 차단할 때 지정하는 값과 원장이 그 기록의
후속 사용을 확인하는 값이 다르면 사용을 막지 못한다.
결국 거래를 승인할 때는 실제 거래의 기록, 감사자가 조사할
기록, 원장이 사용을 거부할 기록이 서로 대응하는지도
검증해야 한다.

본 연구는 이러한 대응을 강제하도록 비공개 거래 검증과
승인된 추적·차단을 연결하는 DPP 프로토콜 zkDPP를 제안한다.
zkDPP는 거래에서 실제 사용·생성한 기록을 기준으로 연결 정보와
사용·차단 식별자를 만들고 암호화했는지를 거래 규칙과 함께
검증한다. 여기서 식별자는 원장이 기록의 사용 여부와 차단
상태를 확인할 때 쓰는 값이다. 이에 따라 일반 참여자는
숨겨진 관계를 열지 않고도 감사 자료가 해당 거래의 기록에
대응함을 확인할 수 있다. 승인된 감사자는 검증된 자료를
복원하여 후속 기록과 남은 차단 대상을 찾고, 원장은 같은
식별자를 기준으로 그 기록을 사용하는 새 거래의 승인을 거부한다.

이 검증에 필요한 조건은 각각의 증명으로 제출하여 서로
연결할 수도 있고, 한 증명에 함께 담을 수도 있다.
여러 증명을 제출하는 구성에서는 그 증명들을 전송·보관하고
반복해서 검증해야 한다. zkDPP는 이러한 부담을 줄이기 위해
계산·이전 규칙과 감사 자료의 대응 조건을 거래 유형별로 묶어,
거래 한 건당 하나의 증명을 제출하도록 구성한다.
이 선택의 비용 효과를 살펴보기 위해 Sahai 등의 두 입력 병합
거래의 증명 분해 방식~\cite{sahai-supply-chain-2020}을 참고한
여섯 증명 구성과 단일 증명 구성을 예비 비교했다.
제한된 비교에서 증명 자체의 크기 합계와 로컬 순차 검증 시간이
감소했지만, 같은 보안 조건의 전체 프로토콜 비교나 현행 zkDPP의
성능 측정은 아니다.~\cite{zkdpp-sahai-microbenchmark-2026}
구체적인 수치와 구현 차이는 \ref{sec:evaluation}절에서 설명한다.

본 논문의 기여는 다음 네 가지다.
\begin{itemize}
\item \textbf{실제 거래에 결합된 감사·차단 구성.}
실제 사용·생성한 기록과 암호화된 연결 정보·사용 제한 식별자의
대응을 검증하여, 거래 검증·감사·차단이 같은 기록을 기준으로
작동하도록 구성한다.
\item \textbf{거래당 하나의 증명 구성.}
거래 유형별 계산·이전 규칙과 감사 자료의 대응 조건을
하나의 증명 관계로 구성한다.
\item \textbf{추적·차단을 포함한 보안 분석.}
명시한 가정 아래 거래·발급·추적·차단의 정확성·건전성과
수동 외부 관찰자의 프라이버시를 분석한다. 승인 이력에 맞는
추적과, 전체 대상의 차단을 확인한 경우에만 완료를 보고하는
성질을 증명한다.
\item \textbf{기능 시험과 비용 평가.}
기호 모델에서 선택한 사례를 검사하고 보존된 시험용 검증기
판본에서 대표 추적·차단 실행을 확인한다.
별도의 성능 측정으로 아홉 거래 유형의 비용과 추적·동결
규모별 결과를 제시한다.~\cite{zkdpp-symbolic-model-2026,
zkdpp-prototype-tests-2026,zkdpp-performance-report-2026}
\end{itemize}

여기서 차단 완료는 조사한 전체 대상이 같은 원장 시점에
아직 사용되지 않았고 차단된 상태임을 확인했다는 뜻이다.
경합이나 자료 부족으로 확인할 수 없으면 미완료로 남기며,
이미 발급된 제품 주장의 취소와 물리적 회수는 포함하지 않는다.
프라이버시 분석은 정상 이력을 읽는 수동 외부 관찰자를 대상으로
하고, 공개 확인 자료와 감사키 복원에 필요한 수보다 적은
키 조각을 가진 경우를 포함한다. 각 분석의 가정과 기능·비용
평가의 판본 및 한계는 본문에서 구분한다.

\newpage
\subsection{관련 연구와 본 연구의 위치}

표~\ref{tab:related-position}는 선행연구가 이미 다루는 기능과
본 논문이 다루는 범위를 나누어 보여 준다. 표의 차이는 각 연구가
제시한 절차의 범위에 관한 것으로, 그 기술을 확장해 같은 기능을
만들 수 없다는 뜻은 아니다.

\begin{table}[htbp]
\centering
\footnotesize
\setlength{\tabcolsep}{2pt}
\caption{선행연구와 본 논문의 비교 범위}
\label{tab:related-position}
\begin{tabularx}{\columnwidth}{@{}>{\raggedright\arraybackslash}p{18mm}YY@{}}
\toprule
연구 & 이미 다루는 기능 & 본 논문의 초점 \\
\midrule
Udokwu--Craß, Babel~\cite{udokwu-crass-2026,babel-hicss-2025}
& 제품 주장 검증; 증빙을 잇는 기록의 이전·재사용 방지
& 증빙 오류 발견 후 후속 기록의 조사와 사용 제한 \\
Sahai 등~\cite{sahai-supply-chain-2020}
& 공급망 문서의 생성·소비와 오염 관련 이력 추적
& 실제 입력 관계의 평상시 은닉, 승인된 복원과 남은 대상의 차단 \\
GGM, Azeroth~\cite{garman-accountable-privacy-2017,jeong-azeroth-2023}
& 거래에 결합된 감사 정보; GGM의 후속 추적
& DPP 거래·감사 관계·차단 식별자의 일치 \\
CAP~\cite{espresso-cap-spec-2022}
& 비공개 자산의 열람과 소유자 협조 없는 동결
& 후속 객체 전체의 선정과 공통 시점 완료 판정 \\
ERC-20R, ACTMeds~\cite{wang-erc20r-2022,sitaula-actmeds-2026}
& 후속 금액의 동결, 의약품 배치의 리콜 제한
& 숨긴 이력에서 도달한 남은 기록 집합의 차단 \\
ZEXE, VeriZexe~\cite{bowe-zexe-2020,xiong-verizexe-2023}
& 비공개 기록의 생성·소비 규칙을 구성할 기반
& DPP 회계·감사·집행과 완료 조건의 구체화 \\
\bottomrule
\end{tabularx}
\end{table}

DPP 연구에서도 여러 단계의 기록과 사용을 이미 다룬다. Babel 등은
생산 단계의 증빙을 원료 지분의 이전·소비와
연결한다.~\cite{babel-hicss-2025}
Udokwu와 Craß는 재활용 원료 함량 등 제품 주장의 조건을 비공개
자료로 검증한다.~\cite{udokwu-crass-2026}
본 논문이 추가로 묻는 것은 그 근거에 문제가
나중에 발견됐을 때 어떤 후속 기록이 아직 사용 가능하며, 그중
무엇의 다음 사용을 막아야 하는가이다. 공급망 구조에서 가장 가까운
Sahai 등은 문서의 생성·소비와 오염 관련 이력을
추적한다.~\cite{sahai-supply-chain-2020}
따라서 추적 자체가 새로운 기능은 아니다. 다만 그 연구에서는
실제 입력 문서의 목록이 공개된다.~\cite[IV절]{sahai-supply-chain-2020}
본 논문은 평상시에 그 연결을 숨기고 승인된 감사자가 복원하는
경우를 다룬다.

감사 자료를 실제 비공개 거래에 결합하는 방법도 GGM과
Azeroth에 이미 있다.~\cite{garman-accountable-privacy-2017,jeong-azeroth-2023}
CAP는 열람 정책 아래 비공개 자산의 열람과 동결을 함께 제공하며,
기존 자산 기록을 소비해 동결 상태의 새 기록을
만든다.~\cite[3.2.4절]{espresso-cap-spec-2022}
본 논문은 찾아낸 원래 기록의
미래 소비 식별자에 현재 상태를 부여한다. 이 저장 방식의 차이만으로
우위를 주장하지는 않는다. 필요한 열람 정책, DPP 전이 규칙,
동결 뒤 새 기록의 대응과 전체 대상 확인을 CAP에 추가한 구성도
같은 사용 제한 목표를 달성할 수 있다.
ERC-20R의 공개 이력상 후속 금액 동결과 ACTMeds의 배치 리콜도
사후 사용 제한의 선행 사례다.~\cite{wang-erc20r-2022,sitaula-actmeds-2026}
다만 각각 금액과 배치를 다루므로,
본 논문이 비공개 생산 관계에서 찾는 미소비 기록의 집합과
차단 완료 조건은 별도로 정의해야 한다. 특히 ERC-20R의 동결은
한 거래에서 원자적으로 실행된다.~\cite[2.1.2절]{wang-erc20r-2022}
본 논문의 여러 상태 변경에는 그런 원자성이 없으며, 소비 경합이
생겨 전체 차단을 확인하지 못하면 완료를 보고하지 않는다.

ZEXE·VeriZexe는 비공개 기록의 생성·소비 규칙을 구성할
기반을 제공한다.~\cite{bowe-zexe-2020,xiong-verizexe-2023}
이러한 기반에 본 논문과 동일한
회계 규칙, 감사 자료, 상태 검사와 완료 절차를 추가한 조합도
같은 목표를 다룰 수 있다. 이는 분석을 위해 구성한 대안이지
두 선행연구에 그 절차가 그대로 구현되어 있다는 뜻은 아니다.
본 논문은 개별 기능의 최초성이나 전체 시스템의 일반적인 성능
우위를 주장하지 않는다. 앞의 예비 관측도 증명 제출·반복 검증
비용에 한정된다. 기여는 DPP 거래에서 실제 입력·출력,
감사자가 복원할 관계, 후속 사용을 차단할 식별자를 일치시키고,
선택한 시점의 남은 대상 전체가 공통 원장 시점에 차단되었을 때만
완료로 판정하는 구체적 구성과 분석에 있다.

이하 \ref{sec:system}절은 역할·권한·신뢰 가정을 정한다.
\ref{sec:accounting}절은 정상적인 회계와 발급을,
\ref{sec:audit}절은 승인된 추적과 차단을 설명한다.
\ref{sec:analysis}절에서 보안 성질을 분석한 뒤
\ref{sec:evaluation}절의 평가와 \ref{sec:limits}절의 한계를 제시한다.
